\documentclass[10pt,letterpaper]{article}
\usepackage[margin=2.2cm]{geometry}
\usepackage{fontspec}
\setmainfont{Arial}
\usepackage[spanish,es-noshorthands]{babel}
\usepackage{xcolor}
\usepackage{booktabs}
\usepackage{tabularx}
\usepackage{colortbl}
\usepackage{graphicx}
\usepackage{hyperref}
\usepackage{microtype}
\usepackage{setspace}

\definecolor{unabblue}{HTML}{173864}
\definecolor{softgray}{HTML}{F9FAFC}
\definecolor{darkgreen}{HTML}{007A33}
\hypersetup{colorlinks=true,linkcolor=unabblue,urlcolor=unabblue}

\onehalfspacing
\setlength{\parindent}{0pt}
\setlength{\parskip}{0.4em}

\begin{document}
\thispagestyle{empty}

\begin{center}
  \includegraphics[width=2.2cm,height=2.2cm,keepaspectratio]{/home/jaimeberrios/Proyectos/UNAB/Ciclo-6/sDSR43/Plantilla-Documentos-Investigacion/assets/img/logo-portada.png}\\[0.2cm]
  {\large\sffamily\bfseries\color{unabblue}UNIVERSIDAD DR. ANDRÉS BELLO}\\[0.05cm]
  {\normalsize\sffamily\bfseries\color{unabblue}FACULTAD DE TECNOLOGÍA E INNOVACIÓN}\\[0.05cm]
  {\small\sffamily\color{unabblue}CARRERA: INGENIERÍA EN SISTEMAS Y COMPUTACIÓN}\\[0.05cm]
  {\footnotesize\sffamily Diseño de Sistemas de Seguridad de Redes de Datos — Ciclo 02-2026}\\[0.3cm]
  {\large\sffamily\bfseries\color{unabblue}INSTRUMENTO DE RECOLECCIÓN DE INFORMACIÓN TÉCNICA}\\[0.1cm]
  {\normalsize\sffamily\bfseries\color{unabblue}Cuestionario de Diagnóstico de Infraestructura y Seguridad de Red}\\[0.15cm]
  {\color{unabblue}\rule{0.6\textwidth}{1.0pt}}
\end{center}

\vspace{0.2cm}

\begin{table}[h!]
  \centering
  \footnotesize
  \begin{tabularx}{\textwidth}{@{}>{\bfseries\color{unabblue}}p{4.2cm} X@{}}
    \toprule
    Unidad de análisis: & Red de datos de la Sede de Enlaces con la Educación (Santa Rosa de Lima). \\
    \midrule
    Informantes clave: & Personal Técnico y Encargada de la Sede Institucional. \\
    \midrule
    Técnica metodológica: & Entrevista técnica estructurada y levantamiento in situ (Hernández-Sampieri, 2023). \\
    \midrule
    Objetivo del instrumento: & Obtener datos sobre topología, obsolescencia del switch central, ausencia de segmentación lógica por VLANs y carencia de controles de filtrado perimetral. \\
    \midrule
    Equipo investigador: & Genesis Saravia, Jaime Berrios, Jefferson Bonilla, Rosmeri González, Florentín Treminio. \\
    \bottomrule
  \end{tabularx}
\end{table}

\vspace{0.2cm}

{\small\itshape\color{black!80}Indicaciones: El presente instrumento tiene como finalidad recabar información operativa y técnica directa de la infraestructura de comunicaciones de la sede. Sus resultados constituyen la base fáctica para el diagnóstico situacional, la matriz de riesgos y la formulación de la propuesta tecnológica.}

\vspace{0.3cm}

\subsection*{\color{unabblue}\small DIMENSIÓN 1: INFRAESTRUCTURA FÍSICA Y EQUIPAMIENTO CENTRAL}

\begin{table}[h!]
  \centering
  \scriptsize
  \renewcommand{\arraystretch}{1.2}
  \begin{tabularx}{\textwidth}{@{}>{\bfseries\color{unabblue}}p{3.2cm} X@{}}
    \toprule
    Ítem 1.1 — Pregunta: & ¿Cuántas estaciones de trabajo operan actualmente en la sede y bajo qué modalidades se distribuye la conectividad hacia los puestos del personal? \\
    \midrule
    Hallazgo técnico: & En la sede operan activamente aproximadamente 10 equipos de cómputo para uso de oficina (1 por cada técnico). La conectividad se distribuye de manera heterogénea: algunos equipos se enlazan mediante cable estructurado hacia el switch o routers intermedios, mientras que otros acceden exclusivamente mediante enlaces Wi-Fi. \\
    Impacto investigativo: & Confirma el tamaño de la unidad de análisis (10 nodos) y la dispersión en los medios de acceso físico y lógico. \\
    \bottomrule
  \end{tabularx}
\end{table}

\begin{table}[h!]
  \centering
  \scriptsize
  \renewcommand{\arraystretch}{1.2}
  \begin{tabularx}{\textwidth}{@{}>{\bfseries\color{unabblue}}p{3.2cm} X@{}}
    \toprule
    Ítem 1.2 — Pregunta: & ¿Cuál es el dispositivo principal que concentra el tráfico del edificio y cuál es su tiempo estimado de operación continua? \\
    \midrule
    Hallazgo técnico: & La distribución central converge en un switch de red instalado desde la construcción de las instalaciones en el año 2003 (aprox. 23 años de servicio). Debido a su obsolescencia, presenta inestabilidad y múltiples puertos físicos dañados que han dejado de transmitir señal. \\
    Impacto investigativo: & Evidencia obsolescencia crítica de hardware como activo principal, sustentando el riesgo extremo en la matriz. \\
    \bottomrule
  \end{tabularx}
\end{table}

\begin{table}[h!]
  \centering
  \scriptsize
  \renewcommand{\arraystretch}{1.2}
  \begin{tabularx}{\textwidth}{@{}>{\bfseries\color{unabblue}}p{3.2cm} X@{}}
    \toprule
    Ítem 1.3 — Pregunta: & ¿Cómo se administran las conexiones salientes del módem del proveedor de servicios de Internet (Claro)? \\
    \midrule
    Hallazgo técnico: & El módem de Claro cuenta con únicamente tres salidas físicas activas: una para el sistema de cámaras de seguridad, otra conectada a un router secundario y la tercera a la computadora principal de la encargada de la sede. Ante la saturación, recientemente se añadió un router adicional. \\
    Impacto investigativo: & Demuestra insuficiencia de puertos y proliferación de routers en cascada sin gestión unificada ni segmentación. \\
    \bottomrule
  \end{tabularx}
\end{table}

\subsection*{\color{unabblue}\small DIMENSIÓN 2: TOPOLOGÍA Y SEGMENTACIÓN LÓGICA DE RED}

\begin{table}[h!]
  \centering
  \scriptsize
  \renewcommand{\arraystretch}{1.2}
  \begin{tabularx}{\textwidth}{@{}>{\bfseries\color{unabblue}}p{3.2cm} X@{}}
    \toprule
    Ítem 2.1 — Pregunta: & ¿Qué topología física rige la comunicación interna y existe actualmente alguna división lógica (VLAN) para separar los flujos de tráfico? \\
    \midrule
    Hallazgo técnico: & La red mantiene una topología en estrella concentrada en el switch. A nivel lógico la arquitectura es totalmente plana: no existe segmentación por VLANs ni delimitación de subredes. Todos los equipos comparten un único dominio de difusión (broadcast) en el mismo rango IP. \\
    Impacto investigativo: & Justifica metodológicamente la necesidad de diseñar un modelo de segmentación lógica por roles operativos. \\
    \bottomrule
  \end{tabularx}
\end{table}

\begin{table}[h!]
  \centering
  \scriptsize
  \renewcommand{\arraystretch}{1.2}
  \begin{tabularx}{\textwidth}{@{}>{\bfseries\color{unabblue}}p{3.2cm} X@{}}
    \toprule
    Ítem 2.2 — Pregunta: & ¿Cuál es el tratamiento y canal de red utilizado cuando ingresan laptops externas de estudiantes o docentes para soporte técnico? \\
    \midrule
    Hallazgo técnico: & Los equipos que ingresan para mantenimiento o diagnóstico se conectan a los mismos puntos de red o redes Wi-Fi compartidas con las estaciones administrativas. No se cuenta con una zona desmilitarizada (DMZ) ni con una subred aislada para cuarentena. \\
    Impacto investigativo: & Valida la existencia de vectores de amenaza de movimiento lateral e infección por malware hacia registros sensibles. \\
    \bottomrule
  \end{tabularx}
\end{table}

\subsection*{\color{unabblue}\small DIMENSIÓN 3: SEGURIDAD PERIMETRAL Y CONTROL DE ACCESO}

\begin{table}[h!]
  \centering
  \scriptsize
  \renewcommand{\arraystretch}{1.2}
  \begin{tabularx}{\textwidth}{@{}>{\bfseries\color{unabblue}}p{3.2cm} X@{}}
    \toprule
    Ítem 3.1 — Pregunta: & ¿Existe actualmente algún servidor central o proxy activo encargado de la auditoría y el filtrado del tráfico hacia Internet? \\
    \midrule
    Hallazgo técnico: & Actualmente no se dispone de ningún servidor de control. En años anteriores operaba un servidor central y un proxy físico dedicado al filtrado del tráfico. No obstante, por modificaciones edilicias fueron retirados, quedando únicamente el switch y el cableado estructurado. \\
    Impacto investigativo: & Fundamenta la propuesta tecnológica de reincorporar un servidor proxy SQUID y cortafuegos perimetral GNU/Linux. \\
    \bottomrule
  \end{tabularx}
\end{table}

\begin{table}[h!]
  \centering
  \scriptsize
  \renewcommand{\arraystretch}{1.2}
  \begin{tabularx}{\textwidth}{@{}>{\bfseries\color{unabblue}}p{3.2cm} X@{}}
    \toprule
    Ítem 3.2 — Pregunta: & ¿Se aplican listas de control de acceso (ACL), políticas de puertos o inspección criptográfica (HTTPS/TLS)? \\
    \midrule
    Hallazgo técnico: & No se cuenta con listas de control de acceso configuradas ni políticas formales para restringir servicios o puertos vulnerables. El tráfico no atraviesa mecanismos de inspección profunda ni políticas forzadas de certificados digitales o filtrado de contenido. \\
    Impacto investigativo: & Establece los requerimientos de diseño para las reglas de filtrado perimetral y políticas de seguridad informática (PSI). \\
    \bottomrule
  \end{tabularx}
\end{table}

\vspace{0.2cm}

{\scriptsize\itshape\color{black!70}Certificación metodológica: Los hallazgos recabados a través del presente instrumento reflejan fielmente el estado tecnológico y de seguridad perimetral de la Sede de Enlaces con la Educación al momento del levantamiento (periodo 2026), sirviendo de insumo directo para la matriz de riesgos y el modelo simulado en GNS3.}

\begin{center}
  \vspace{0.1cm}
  {\footnotesize\sffamily\bfseries\color{unabblue}San Miguel, El Salvador — Equipo de Investigación de Cátedra UNAB}
\end{center}

\end{document}
